%%
%% Starter template for SOSP submissions.
%%
%% Source: https://sigops.org/s/conferences/sosp/2026/cfp.html
%%
%% Preamble is as specified by the SOSP 2026 CFP.
%% Always verify against the current year's CFP before submitting.
%%

\documentclass[sigplan,10pt]{acmart}
\renewcommand\footnotetextcopyrightpermission[1]{}   % suppress copyright block
\settopmatter{printfolios=true}                       % show page numbers

%% ============================================================
%% Packages
%% ============================================================
\usepackage{booktabs}    % professional tables (\toprule, \midrule, \bottomrule)
\usepackage{graphicx}    % \includegraphics
\usepackage{subcaption}  % subfigures
\usepackage{microtype}   % better text justification
\usepackage{xspace}      % smart spacing after macros
%% Placeholder tracking — copy utils/fillin.sty to this directory.
%% Use \fillin[desc], \fillcite[desc], \fillnum for unfinished values.
%% Add \printfillinlistifany before \end{document} during drafting.
%% Remove this line before submission.
\usepackage{fillin}

%% ============================================================
%% Custom macros
%% ============================================================
\newcommand{\sysname}{SysName\xspace}   % replace with your system's name
\newcommand{\eg}{e.g.,\xspace}
\newcommand{\ie}{i.e.,\xspace}
\newcommand{\etal}{\textit{et al.}\xspace}

%% ============================================================
\begin{document}

\title{\sysname: A [Noun Phrase Capturing Your Key Contribution]}

%% For blind submission: leave author blocks in but anonymized,
%% or follow the current CFP guidance on anonymization.
\author{Author One}
\affiliation{\institution{Anonymous Institution}\country{}}
\author{Author Two}
\affiliation{\institution{Anonymous Institution}\country{}}

\begin{abstract}
%% 5-sentence formula:
%% 1. The problem (concrete, with a number if possible)
%% 2. Why existing systems fail to solve it
%% 3. Key insight + high-level design approach
%% 4. Evaluation setup: workload(s), baseline(s), scale
%% 5. Headline result

[1. Concrete problem statement.] [2. Why existing systems fail.] [3. Key insight:
we observe that X, which enables our design to do Y.] [4. We implement \sysname
in N~kLoC of [language] and evaluate it on [workload] against [baseline] on a
[N-node cluster].] [5. \sysname achieves [Nx] improvement over [baseline] on
[benchmark], while reducing [overhead] by [Y\%].]
\end{abstract}

\maketitle
\pagestyle{plain}   % required by SOSP CFP

%% ============================================================
\section{Introduction}
%% ============================================================

%% Opening scenario: concrete situation where existing systems fail (with numbers).
[Opening scenario.]

%% Fundamental challenge: why existing approaches are structurally insufficient.
[The challenge.]

%% State the key insight explicitly.
\textbf{Key insight.}
We observe that [one-sentence insight]. This enables [high-level design approach].

%% Brief design overview.
We design \sysname, [one-sentence description]. \S\ref{sec:design} describes the
design in detail.

%% Contribution bullets (3–5, specific and falsifiable).
This paper makes the following contributions:
\begin{itemize}
  \item We identify [observation] and show that [evidence] (\S\ref{sec:motivation}).
  \item We design \sysname, achieving [goal] via [mechanism] (\S\ref{sec:design}).
  \item We implement \sysname in [N]k~LoC of [language] (\S\ref{sec:impl}).
  \item We show [Nx] improvement over [baseline] on [benchmark] (\S\ref{sec:eval}).
\end{itemize}

\S\ref{sec:motivation} presents a motivating study.
\S\ref{sec:design} describes the design.
\S\ref{sec:impl} covers implementation.
\S\ref{sec:eval} presents evaluation.
\S\ref{sec:related} discusses related work.
\S\ref{sec:conclusion} concludes.


%% ============================================================
\section{Motivation}
\label{sec:motivation}
%% ============================================================
%% Must: (1) measure existing systems failing on a real workload,
%%       (2) explain WHY they fail structurally,
%%       (3) connect the failure to your key insight.

\subsection{The Problem in Practice}

[Concrete scenario. Figure~\ref{fig:motivation} shows the measured failure.]

\begin{figure}[t]
  \centering
  \includegraphics[width=\columnwidth]{figs/motivation.pdf}
  \caption{[Existing system] under [workload]: [metric] rises from [A] to [B]
           when [condition], violating the [X~ms] SLO for [Z\%] of requests.}
  \label{fig:motivation}
\end{figure}

\subsection{Why Existing Systems Fall Short}

[Structural explanation—not just "they are slow" but WHY their design cannot
avoid this problem.]

\subsection{Key Insight}

[Restate the insight. Connect it to the failure mode above.]


%% ============================================================
\section{Design}
\label{sec:design}
%% ============================================================

\sysname is designed to achieve:
\begin{enumerate}
  \item \textbf{Goal 1.} [Specific, measurable goal.]
  \item \textbf{Goal 2.} [Specific, measurable goal.]
\end{enumerate}
\sysname does \emph{not} aim to [non-goal].

\subsection{Overview}

[High-level description with Figure~\ref{fig:arch}.]

\begin{figure}[t]
  \centering
  \includegraphics[width=\columnwidth]{figs/architecture.pdf}
  \caption{\sysname architecture. [Component A] handles [X]; [Component B]
           manages [Y]. Shaded = new; unshaded = existing infrastructure.}
  \label{fig:arch}
\end{figure}

\subsection{Component A}
[Detail, interface, invariants, key trade-offs.]

\subsection{Component B}
[Detail.]

\subsection{Correctness}
[Argue key safety/liveness properties, or discuss failure handling.]


%% ============================================================
\section{Implementation}
\label{sec:impl}
%% ============================================================

We implement \sysname in [language], comprising approximately [N]k lines of code,
running on [OS, kernel version].

\paragraph{[Key aspect 1.]} [Description.]
\paragraph{[Key aspect 2.]} [Description.]


%% ============================================================
\section{Evaluation}
\label{sec:eval}
%% ============================================================
%% Every claim from §1 bullets must have a corresponding experiment.

\subsection{Experimental Setup}

\paragraph{Testbed.}
[CPU model, cores, NUMA; Memory; Storage (NVMe/HDD); Network; OS + kernel.]

\paragraph{Baselines.}
\begin{itemize}
  \item \textbf{[Baseline 1]} [version]: [why it is the right comparison.]
  \item \textbf{[Baseline 2]} [version]: [description.]
\end{itemize}

\paragraph{Workloads.}
[Workload name, configuration, and parameters for each benchmark used.]

\subsection{Main Results}

\textbf{Claim:} \sysname achieves [Nx] higher throughput than [baseline] on [workload].

[Describe Figure~\ref{fig:main}. Tell the reader what to observe and why it
demonstrates the claim.]

\begin{figure}[t]
  \centering
  \includegraphics[width=\columnwidth]{figs/main.pdf}
  \caption{Throughput vs.\ p99 latency on [workload]. \sysname achieves
           [Nx] higher throughput at the [X~ms] p99 target.
           Error bars = one standard deviation over 5 runs.}
  \label{fig:main}
\end{figure}

\subsection{Scalability}
[Throughput/speedup vs.\ number of cores or nodes.]

\subsection{Overhead}
[Overhead on the unmodified workload / common case, as \% above baseline.]

\subsection{Microbenchmarks}
[Component-level benchmarks, each tied back to a specific claim.]


%% ============================================================
\section{Related Work}
\label{sec:related}
%% ============================================================
%% Organize THEMATICALLY, not paper-by-paper.

\paragraph{[Theme 1: e.g., distributed KV stores].}
[Group of papers doing X by Y. How \sysname differs and why.]

\paragraph{[Theme 2].} [...]


%% ============================================================
\section{Conclusion}
\label{sec:conclusion}
%% ============================================================

We present \sysname, [one-sentence summary]. Our key insight is that [restate].
\sysname achieves [headline result] over [baseline] on [workload].


%% Acknowledgments — remove or anonymize for submission
%\begin{acks}
%We thank [...]
%\end{acks}

\bibliographystyle{ACM-Reference-Format}
\bibliography{references}

\end{document}
